
\section{Related Work}
\label{sec:related_work}

\subsection{World Models and Video World Models} 
A world model takes in an observation of a system along with a chosen action or intervention, and produces a prediction of the resulting state or outcome. The notion of a world model can be traced back to Kenneth Craik, who argued in his 1943 book \textit{The Nature of Explanation}~\cite{craik1943nature} that if an organism carries a small internal model of reality and its possible actions, it can test options mentally, anticipate the future, use past experience, and respond more effectively and safely to challenges. These ideas later became central to dynamic systems and control theory~\cite{bellman1957dynamic, kalman1960new, bryson1975applied}, establishing mathematical tools for modeling, prediction, and planning in uncertain environments. The Dyna architecture proposed by Sutton~\cite{sutton_dyna1991} emphasized that learning an internal model of the world enables agents to plan their actions rather than rely solely on reactive trial-and-error interactions with the environment, such as in model-free reinforcement learning. More recently, the concept of world models has been combined with deep generative modeling. Several works ~\cite{ha2018world, hafner2020dreamcontrollearningbehaviors, hafner2020mastering, hafner2023mastering, hafner2025trainingagentsinsidescalable} demonstrated learning a compact latent dynamics model directly from pixel observations, enabling policy optimization within a learned latent environment. Instead of modeling raw pixels or compression-based latent spaces such as those from variational autoencoders~\cite{kingma2013auto} that remain close to pixel space~\cite{zheng2025diffusion}, an important direction is for world models to learn abstract representations where predictions can ignore irrelevant details, enabling better understanding, prediction, and planning~\cite{zhou2024dino, bardes2024revisiting, assran2025v}.

With the emergence of video diffusion transformers~\cite{peebles2023scalablediffusionmodelstransformers} and major advances in text-to-video and image-to-video generation, a growing trend in world modeling is to use video diffusion models as \emph{world simulators}~\cite{openai2024sora,yang2023unisim, liang2024dreamitate, bar2025navigation}. Since this line of work typically does not focus on learning abstract representations or on planning and decision making within them, we refer to these approaches as \emph{video world models} to emphasize their reliance on video diffusion for simulation. Video world models have been applied to numerous domains including robotics~\cite{mereu2025generative, yang2024learninginteractiverealworldsimulators, li2025unified,li2026causal,ye2026world,gao2026dreamdojo}, video games~\cite{oasis2024, parkerholder_fruchter_2025_genie3, valevski2025diffusionmodelsrealtimegame, alonso2024diffusionworldmodelingvisual, he2025matrixgame20opensourcerealtime, xiao2025worldmem}, self-driving~\cite{agarwal2025cosmos,hu2023gaia,waymo2026wm}, and physical simulation~\cite{li2025pisa, yuan2026inference}. Our work contributes to this direction. However, our focus is on building a \emph{multiplayer} video world model. Instead of just generating pixels from actions, we want to create a foundation for agent learning in a setting where multiple agents share the same world, and where simply simulating pixels is not enough. To our knowledge, Multiverse~\cite{enigma2025multiverse} is the only video world model capable of simulating multiple agents. While their U-Net~\cite{ronneberger2015u} model is trained on gameplay from one race track (``Tsukuba'') in the 2004 video game Gran Turismo 4, we tackle the significantly more complex 3D open-world environment of Minecraft. We study applying their channel-concatenation design to our model in~\cref{sec:experiments}.








\subsection{Autoregressive Video Generation}

Large scale Diffusion Transformer (DiT)~\cite{peebles2023scalablediffusionmodelstransformers} video generation models have made tremendous progress in recent years~\cite{openai2024sora, kong2025hunyuanvideosystematicframeworklarge, wan2025wanopenadvancedlargescale}. Diffusion Forcing~\cite{chen2024diffusionforcingnexttokenprediction} is a technique where autoregressive generation emerges as a byproduct of training with an independent noise level per frame. CausVid~\cite{yin2025slowbidirectionalfastautoregressive} was introduced to convert a bidirectional video model to an efficient causal model of comparable quality. Self-Forcing~\cite{huang2025selfforcingbridgingtraintest,cui2025self, hong2025relic} improves generation quality further by supervising on a model's own generations, mitigating autoregression train-test mismatch.  

A concurrent work to ours, RELIC~\cite{hong2025relic}, also studies Self-Forcing with a long-context teacher. Like our method, they propose a memory-efficient implementation involving a recomputation step for the backward pass. However, we distinguish our approach by performing this step in one parallel forward pass via masking, avoiding the multiple rolling passes required by RELIC.

\subsection{AI Agents in Minecraft}

A number of platforms for agents in Minecraft have been developed, which have primarily been used for research in reinforcement learning~\cite{johnson2016malmo, fan2022minedojo, guss2019minerllargescaledatasetminecraft}. Mineflayer~\cite{mineflayer2025} is a commonly used framework for developing bots in Minecraft, and was used to collect the LoopNav~\cite{lian2025toward} dataset testing spatial memory, as well as in Voyager~\cite{wang2023voyager}, a Minecraft LLM agent. VPT~\cite{baker2022videopretrainingvptlearning} is a large-scale dataset of single-player Minecraft data collected from humans. However, as we explain in~\cref{sec:system}, none of these frameworks are capable of simulating multiplayer gameplay with visuals, leading us to develop \solarisengine. 
